\documentclass[10pt,a4paper]{amsart}
\usepackage{amsmath,amssymb,amsthm,graphicx}
\usepackage{verbatim}
\newenvironment{reference}{\comment}{\endcomment}
%\newenvironment{reference}{}{}
\newenvironment{slogan}{\comment}{\endcomment}
\newenvironment{history}{\comment}{\endcomment}

% For commutative diagrams we use Xy-pic
\usepackage[all]{xy}

% We use 2cell for 2-commutative diagrams.
\xyoption{2cell}
% Theorem environments.
%
\theoremstyle{plain}
\newtheorem{theorem}[subsection]{Theorem}
\newtheorem{proposition}[subsection]{Proposition}
\newtheorem{lemma}[subsection]{Lemma}

\theoremstyle{definition}
\newtheorem{definition}[subsection]{Definition}
\newtheorem{example}[subsection]{Example}
\newtheorem{exercise}[subsection]{Exercise}
\newtheorem{situation}[subsection]{Situation}

\theoremstyle{remark}
\newtheorem{remark}[subsection]{Remark}
\newtheorem{remarks}[subsection]{Remarks}

\numberwithin{equation}{subsection}

% Macros
%
\def\lim{\mathop{\mathrm{lim}}\nolimits}
\def\colim{\mathop{\mathrm{colim}}\nolimits}
\def\Spec{\mathop{\mathrm{Spec}}}
\def\Hom{\mathop{\mathrm{Hom}}\nolimits}
\def\Ext{\mathop{\mathrm{Ext}}\nolimits}
\def\SheafHom{\mathop{\mathcal{H}\!\mathit{om}}\nolimits}
\def\SheafExt{\mathop{\mathcal{E}\!\mathit{xt}}\nolimits}
\def\Sch{\mathit{Sch}}
\def\Mor{\mathop{\mathrm{Mor}}\nolimits}
\def\Ob{\mathop{\mathrm{Ob}}\nolimits}
\def\Sh{\mathop{\mathit{Sh}}\nolimits}
\def\NL{\mathop{N\!L}\nolimits}
\def\CH{\mathop{\mathrm{CH}}\nolimits}
\def\proetale{{pro\text{-}\acute{e}tale}}
\def\etale{{\acute{e}tale}}
\def\QCoh{\mathit{QCoh}}
\def\Ker{\mathop{\mathrm{Ker}}}
\def\Im{\mathop{\mathrm{Im}}}
\def\Coker{\mathop{\mathrm{Coker}}}
\def\Coim{\mathop{\mathrm{Coim}}}

% Boxtimes
%
\DeclareMathSymbol{\boxtimes}{\mathbin}{AMSa}{"02}

%
% Macros for moduli stacks/spaces
%
\def\QCohstack{\mathcal{QC}\!\mathit{oh}}
\def\Cohstack{\mathcal{C}\!\mathit{oh}}
\def\Spacesstack{\mathcal{S}\!\mathit{paces}}
\def\Quotfunctor{\mathrm{Quot}}
\def\Hilbfunctor{\mathrm{Hilb}}
\def\Curvesstack{\mathcal{C}\!\mathit{urves}}
\def\Polarizedstack{\mathcal{P}\!\mathit{olarized}}
\def\Complexesstack{\mathcal{C}\!\mathit{omplexes}}
% \Pic is the operator that assigns to X its picard group, usage \Pic(X)
% \Picardstack_{X/B} denotes the Picard stack of X over B
% \Picardfunctor_{X/B} denotes the Picard functor of X over B
\def\Pic{\mathop{\mathrm{Pic}}\nolimits}
\def\Picardstack{\mathcal{P}\!\mathit{ic}}
\def\Picardfunctor{\mathrm{Pic}}
\def\Deformationcategory{\mathcal{D}\!\mathit{ef}}
\usepackage{cleveref}
\crefname{originallemma}{original lemma}{original lemmas}
\Crefname{originallemma}{Original lemma}{Original lemmas}
\makeatletter
\newlabel{categories-lemma-finite-limits-exist}{{002O}{source}{}{sourcepointer-1}{}}
\newlabel{categories-lemma-finite-limits-exist@cref}{{[originallemma][0][]002O}{[1][0][]source}{}{}{}}
\makeatother
\title{A conservative family of points detects injectivity, surjectivity, equality and limits}
\author{The Stacks Project Authors}
\date{stacks-project:sites.tex:local-snapshot-20261012; GNU FDL 1.2 or later}
\begin{document}
\maketitle
\noindent\textit{Editorial note.} All delivered mathematical text is the author's own, copied from the pinned Stacks Project chapter file sites-cohomology.tex; no statement or proof step was rewritten, reordered or removed.\par\medskip
\noindent\textit{Editorial note.} The proof is the author's complete argument; every step the author writes out is inside the counted span and no closing sentence is a placeholder.\par\medskip
\noindent\textit{Editorial note.} Every labelled reference occurring inside the counted spans is carried in assembly.external\_references as a non-load-bearing navigation pointer (same-chapter ones at their definition site, cross-chapter ones at the invocation inside the counted span), so no load-bearing mathematics is replaced by a stub. No bibliography citation occurs inside the counted spans.\par\medskip
\noindent\textit{Editorial note.} The checked chapter list for same-name collisions is the set of stacks chapters staged in this repository: chow, cohomology, divisors, etale-cohomology, flat, more-algebra, more-morphisms, morphisms, spaces-more-morphisms, varieties together with this card's chapter sites-cohomology.\par\medskip
\noindent\textit{Editorial note.} The preamble is the chapter's own retained minimal preamble (sources/stacks-sites-cohomology/retained-preamble.tex), a verbatim excerpt of the native preamble: the Xy-pic setup and the native theorem declarations with the mathematical macros.\par\medskip
\noindent\textit{Allowed standard prerequisites (editorial).}\begin{itemize}
\item Stalk functors commute with finite limits and colimits. Same chapter, the statement that the stalk functors are exact. Allows the conditions to be checked after applying the stalk functors.
\item Injections and surjections as isomorphisms criteria. Same chapter, the characterization of injections by the diagonal and of surjections by the codiagonal being isomorphisms, together with Lemma lemma-mono-epi-sheaves. Turns injectivity and surjectivity into isomorphism statements to which conservativity applies.
\item Conservative families of points. Same chapter, the definition of a conservative family of points. Turns pointwise isomorphism into isomorphism.
\item Equalizers as finite limits. Categories, Lemma categories-lemma-finite-limits-exist. Expresses the equalizer in terms of products and fibre products.
\end{itemize}\medskip
\ifdefined\hypertarget\hypertarget{sourcepointer-1}{}\fi
\noindent\textit{Source reference.} Original lemma 002O: sources/stacks-sites/pinned-original-source.tex, source lines 9304--9304. Navigation pointer invoked at proof line 9304; the counted span contains the pointer only and the cited result is not reproduced inside it, so the pointer is registered as a non-load-bearing navigation stub in the established form for Stacks material. No bibliography citation is involved.\par\medskip
\begin{lemma}
\label{lemma-exactness-stalks}
Let $\mathcal{C}$ be a site and let $\{p_i\}_{i\in I}$ be a conservative
family of points. Then
\begin{enumerate}
\item Given any map of sheaves $\varphi : \mathcal{F} \to \mathcal{G}$
we have $\forall i, \varphi_{p_i}$ injective implies $\varphi$ injective.
\item Given any map of sheaves $\varphi : \mathcal{F} \to \mathcal{G}$
we have $\forall i, \varphi_{p_i}$ surjective implies $\varphi$ surjective.
\item Given any pair of maps of sheaves
$\varphi_1, \varphi_2 : \mathcal{F} \to \mathcal{G}$
we have $\forall i, \varphi_{1, p_i} = \varphi_{2, p_i}$
implies $\varphi_1 = \varphi_2$.
\item Given a finite diagram $\mathcal{G} : \mathcal{J}
\to \Sh(\mathcal{C})$, a sheaf $\mathcal{F}$ and morphisms
$q_j : \mathcal{F} \to \mathcal{G}_j$ then $(\mathcal{F}, q_j)$
is a limit of the diagram if and only if for each $i$ the stalk
$(\mathcal{F}_{p_i}, (q_j)_{p_i})$ is one.
\item Given a finite diagram $\mathcal{F} : \mathcal{J}
\to \Sh(\mathcal{C})$, a sheaf $\mathcal{G}$ and morphisms
$e_j : \mathcal{F}_j \to \mathcal{G}$ then $(\mathcal{G}, e_j)$
is a colimit of the diagram if and only if for each $i$ the stalk
$(\mathcal{G}_{p_i}, (e_j)_{p_i})$ is one.
\end{enumerate}
\end{lemma}

\begin{proof}
We will use over and over again that all the stalk functors commute
with any finite limits and colimits and hence with products, fibred
products, etc. We will also use that injective maps are the monomorphisms
and the surjective maps are the epimorphisms.
A map of sheaves $\varphi : \mathcal{F} \to \mathcal{G}$
is injective if and only if
$\mathcal{F} \to \mathcal{F} \times_\mathcal{G} \mathcal{F}$
is an isomorphism. Hence (1).
Similarly, $\varphi : \mathcal{F} \to \mathcal{G}$
is surjective if and only if
$\mathcal{G} \amalg_\mathcal{F} \mathcal{G} \to \mathcal{G}$
is an isomorphism. Hence (2).
The maps $a, b : \mathcal{F} \to \mathcal{G}$
are equal if and only if
$\text{Equalizer}(a, b : \mathcal{F} \to \mathcal{G}) \to \mathcal{F}$
is an isomorphism\footnote{The equalizer can be expressed
in terms of products and fibre products, see proof of
Categories, Lemma \ref{categories-lemma-finite-limits-exist}.}. Hence (3).
The assertions (4) and (5) follow immediately from the definitions
and the remarks at the start of this proof.
\end{proof}

\end{document}
